\section{Three-body verification of the transfer}
\label{sec:cr3bp}

\subsection{Motivation}

Section~\ref{sec:perturbations} treated the Earth as a perturbation on a
Moon-centred two-body problem. This part removes that hierarchy and re-solves
the transfer with both primaries on an equal footing, in the Earth--Moon
circular restricted three-body problem. The purpose is not a prettier
trajectory -- it is visually identical -- but a bound on the modelling error of
the two-body design.

\subsection{Formulation}

In the synodic frame, rotating with the primaries at unit rate and scaled so
that the Earth sits at $x = -\mu$ and the Moon at $x = 1-\mu$ with
$\mu = 0.01215$, the equations of motion are \cite{szebehely}

\begin{align}
\ddot x &= 2\dot y + x - \frac{(1-\mu)(x+\mu)}{r_1^3} - \frac{\mu(x-1+\mu)}{r_2^3},\label{eq:cr3bpx}\\
\ddot y &= -2\dot x + y - \frac{(1-\mu)y}{r_1^3} - \frac{\mu y}{r_2^3},\label{eq:cr3bpy}\\
\ddot z &= -\frac{(1-\mu)z}{r_1^3} - \frac{\mu z}{r_2^3},\label{eq:cr3bpz}
\end{align}

with $r_1$ and $r_2$ the distances to Earth and Moon. Length is scaled by
$LU = d_{EM}$ and time by $TU = \sqrt{LU^3/(\mu_E+\mu_M)} = 375\,191\unit{s}$,
giving a velocity unit $VU = 1.0245\unit{km/s}$.

\subsection{Frame conversion}

Converting a Moon-centred inertial state into this frame takes three steps, in
order: rotate the components into the frame whose $+x$ is the Earth-to-Moon
line; translate the origin from the Moon to the barycentre and rescale; then
subtract the frame rotation and add the Moon's own inertial velocity, which in
these units is $[0, 1-\mu, 0]^\top$. Writing $\alpha$ for the inertial longitude
of the Earth-to-Moon direction,

\begin{equation}
\mathbf r_{\text{syn}} = R_z(-\alpha)\frac{\mathbf r_{MCI}}{LU} + \begin{bmatrix}1-\mu\\0\\0\end{bmatrix},
\qquad
\mathbf v_{\text{syn}} = R_z(-\alpha)\frac{\mathbf v_{MCI}}{VU} + \begin{bmatrix}0\\1-\mu\\0\end{bmatrix}
                        - \hat{\mathbf z}\times\mathbf r_{\text{syn}} .
\label{eq:syn}
\end{equation}

The $(1-\mu)$ that looks mysterious in component form is simply the Moon's
velocity about the barycentre. The conversion is unit-tested by an
unambiguous case: a body at rest at the Moon's centre must have exactly zero
synodic velocity, because in the rotating frame the Moon does not move. The
implementation returns zero to machine precision, and the Earth lands at
$x = -\mu$ to $1.7\times10^{-16}$.

Consistency with Section~\ref{sec:perturbations} is maintained by deriving
$\alpha$ from the same Earth ephemeris used there, so the two perturbation
studies describe the same geometry rather than two different ones.

\subsection{Single shooting}

The decision vector is $\mathbf p = [\delta v_x, \delta v_y, t_f]$ in synodic
non-dimensional units, initialised from the two-body design of
Section~\ref{sec:transfer}. Both vehicles are propagated with
\eqref{eq:cr3bpx}--\eqref{eq:cr3bpz} and the cost is the Euclidean miss at $t_f$,
minimised with Nelder--Mead.

One subtlety deserves a note. The problem is planar, so the miss has only two
independent components against three decision variables: the zero-miss set is a
\emph{curve}, not a point. An unregularised optimiser wanders along that curve
and reports an arbitrary member of a family, each with a slightly different
$\dv$ and time of flight. A small penalty on the distance from the two-body
guess is therefore added, weighted so that it can never dominate the miss term
-- a $1\%$ deviation costs one centimetre of apparent miss. The reported
correction is then the \emph{smallest} one on that curve, which is the only
version of the number that means anything.

\subsection{Results}

\begin{table}[htbp]
\centering
\caption{Two-body design against the CR3BP solution. The correction is
millimetres per second on a $118.80\unit{m/s}$ budget: a relative effect of
order $10^{-4}$.}
\label{tab:cr3bp}
\small
\begin{tabular}{lrrr}
\toprule
Quantity & Two-body & CR3BP & $|\Delta|$ \\
\midrule
$\dv_1$ [km/s]  & $0.060524$ & $0.060522$ & $0.000002$ \\
$\dv_2$ [km/s]  & $0.058276$ & $0.058269$ & $0.000007$ \\
$\dv_{\text{tot}}$ [km/s] & $0.118800$ & $0.118791$ & $\mCRextra\times10^{-6}$ \\
Time of flight [s] & \mTOF & \mCRtof & $\mCRtofShift$ \\
\bottomrule
\end{tabular}
\end{table}

The shooter converges to a miss of $\mCRmiss\unit{m}$, and the total correction
to the two-body design is $\mCRextra\unit{mm/s}$ with a time-of-flight shift of
$\mCRtofShift\unit{s}$, i.e.\ $0.02\%$ of the coast.

The operational conclusion is worth stating plainly, because it justifies the
whole two-body approach of the preceding sections: a $100$--$400\unit{km}$ lunar
orbit sits deep enough in the lunar potential well that including the Earth as a
full primary changes the impulsive budget by millimetres per second. That is
below the noise of any realistic navigation solution. The three-body model is
the right tool for a translunar cruise or a halo orbit; for this mission it
confirms rather than corrects.

\begin{figure}[htbp]
\centering
\includegraphics[width=\textwidth]{fig16_cr3bp_transfer.png}
\caption{Optimised transfer in the Earth--Moon synodic frame, Moon-centred view.
The arc is the same climb as Figure~\ref{fig:geometry} seen from a rotating
frame. Right: range history and the correction relative to the two-body design,
in millimetres per second.}
\label{fig:cr3bp}
\end{figure}

\subsection{Verification}

The CR3BP admits no energy integral, but the Jacobi constant

\begin{equation}
C_J = 2\left[\frac{x^2+y^2}{2} + \frac{1-\mu}{r_1} + \frac{\mu}{r_2}\right] - v^2
\label{eq:jacobi}
\end{equation}

is exactly conserved along any ballistic arc. Evaluated along the optimised
transfer it drifts by $9\times10^{-13}$ relative, a pure measure of integration
quality that is independent of every other check in this report.

The frame conversion is verified by the three cases listed above, and the
round trip inertial $\to$ synodic $\to$ inertial closes to $10^{-11}$.
